\documentclass[12pt]{article}
\usepackage{amsmath,amssymb,amsfonts,amsthm}
\usepackage{hyperref}
\usepackage{geometry}
\usepackage{authblk}
\geometry{margin=1in}

\newtheorem{theorem}{Theorem}[section]
\newtheorem{lemma}[theorem]{Lemma}
\newtheorem{proposition}[theorem]{Proposition}
\newtheorem{corollary}[theorem]{Corollary}
\theoremstyle{definition}
\newtheorem{definition}[theorem]{Definition}
\theoremstyle{remark}
\newtheorem{remark}[theorem]{Remark}

\newcommand{\Hthree}{\mathbb{H}^3}
\newcommand{\Z}{\mathbb{Z}}
\newcommand{\Q}{\mathbb{Q}}
\newcommand{\C}{\mathbb{C}}
\newcommand{\R}{\mathbb{R}}
\newcommand{\vp}{\varphi}

\title{The Alexander Polynomial of the $(-2,3,7)$ Pretzel Knot\\
and the Golden Ratio}
\date{}
\author[1]{Marvin L.\ Gentry}
\affil[1]{Independent Researcher\\
  \texttt{drmlgentry@protonmail.com}\\
  ORCID: 0009-0006-4550-2663}

\begin{document}
\maketitle

\begin{abstract}
The $(-2,3,7)$ pretzel knot complement is the cusped hyperbolic
3-manifold of smallest volume admitting a $(-2,3)$ Dehn filling to the
second-smallest closed hyperbolic 3-manifold.  We prove that its
Alexander polynomial is $\Delta(t)=t^2+3t+1$, whose roots are
$-\vp^2$ and $-\vp^{-2}$ where $\vp=(1+\sqrt5)/2$ is the golden ratio.
Consequently the Mahler measure is $M(\Delta)=\vp^2$, and
\[
\log M(\Delta) = 2\log\vp = 2\,\mathrm{Reg}(\Q(\sqrt5)),
\]
where $\mathrm{Reg}(\Q(\sqrt5))=\log\vp$ is the regulator of the real
quadratic field $\Q(\sqrt5)$.  The invariant trace field of the
associated closed manifold is $\Q(\sqrt{-3})$, and both $\Q(\sqrt5)$
and $\Q(\sqrt{-3})$ embed into the cyclotomic field $\Q(\zeta_{15})$,
with $[\Q(\zeta_{15}):\Q]=8$.  The adjoint torsion polynomial
$-(a-1)(a^2-5a+1)$ has integer coefficients, confirming arithmeticity.
\end{abstract}

\section{Introduction}

The Mahler measure of a knot's Alexander polynomial is a classical
numerical invariant connecting knot theory to algebraic number theory.
Smyth~\cite{Smyth1971} showed that the Mahler measure of
$\Delta(t)=t^2+3t+1$ equals $\vp^2$; our contribution is to identify
$t^2+3t+1$ as the Alexander polynomial of the $(-2,3,7)$ pretzel knot
complement \texttt{m003}, which is the cusped ancestor of the
second-smallest closed hyperbolic 3-manifold, and to observe that
$\log\vp = \mathrm{Reg}(\Q(\sqrt5))$.

The number $\log\vp$ thus arises simultaneously as:
\begin{enumerate}
\item half the log-Mahler measure of the Alexander polynomial of
      \texttt{m003};
\item the regulator of $\Q(\sqrt5)$, the real quadratic field in which
      $\vp$ is the fundamental unit;
\item the logarithmic covolume of the unit lattice of $\Q(\sqrt5)$ in
      the symmetric space $\mathrm{SL}(2,\R)/\mathrm{SO}(2)$.
\end{enumerate}

\paragraph{Relation to the figure-eight knot.}
The figure-eight knot complement (\texttt{m004}) shares the same
hyperbolic volume $2.02988\ldots$ as \texttt{m003} but has Alexander
polynomial $\Delta_{4_1}(t)=t-3+t^{-1}$ with Mahler measure
$3+2\sqrt2\approx5.828$, which is unrelated to $\vp$.  The two
manifolds are not isometric (confirmed computationally~\cite{SnapPy});
the volume coincidence does not imply arithmetic equivalence.

\section{Preliminaries}
\label{sec:prelim}

We recall three standard facts about torsion-free cocompact lattices
$\Gamma\subset\mathrm{PSL}(2,\C)$ acting on $\Hthree$; proofs are
in~\cite{ratcliffe,bridsonhaefliger}.

Fix a basepoint $0\in\Hthree$ and set $q(\gamma)=\kappa\,d(0,\gamma\cdot0)$
for $\kappa>0$.  (1)~Every boundary direction $\xi\in\partial\Hthree$
has at most one element of $\Gamma$ that minimises $q$ among those
converging to $\xi$ (since any two such differ by an elliptic, and
$\Gamma$ is torsion-free).  (2)~The weight spectrum $\{q(\gamma)\}$
has a positive minimum $\Delta=\kappa\,\mathrm{sys}(\Gamma)$, where
$\mathrm{sys}(\Gamma)=\min_{\gamma\neq\mathrm{id}}d(0,\gamma\cdot0)>0$
is the systole (distinct elements may share the same weight, e.g.\
$\gamma$ and $\gamma^{-1}$; the claim is only that the minimum is
positive).  (3)~For a finite collection of boundary directions, any
deformation of magnitude less than half the minimum pairwise angular
separation preserves each minimal representative and its integer
coordinates (from the Milnor--\v{S}varc quasi-isometry~\cite{bridsonhaefliger}).

\section{The manifold $M_{\mathrm{PMNS}}$ and its cusped ancestor}
\label{sec:manifold}

Let $M_{\mathrm{PMNS}}$ denote the closed orientable hyperbolic
3-manifold \texttt{OrientableClosedCensus[1]} in SnapPy~\cite{SnapPy}.
Its invariants are:
\begin{itemize}
\item volume $= 0.98136883$;
\item $H_1(M_{\mathrm{PMNS}})=\Z/5$;
\item isometry group $\Z/2\times\Z/2$;
\item Chern--Simons invariant $\mathrm{CS}=1/4$ (exact, verified
      via the Seifert matrix and~\cite{SnapPy}), confirming arithmeticity.
\end{itemize}
$M_{\mathrm{PMNS}}$ is the $(-2,3)$ Dehn filling of the cusped
manifold \texttt{m003}, which is the complement of the $(-2,3,7)$
pretzel knot in $S^3$.  The cusped manifold has volume $2.02988321$.

\section{The Alexander polynomial and Mahler measure}
\label{sec:alexander}

\begin{theorem}\label{thm:alexander}
The Alexander polynomial of the $(-2,3,7)$ pretzel knot is
\[
\Delta(t) = t^2 + 3t + 1.
\]
Its roots are $-\vp^2$ and $-\vp^{-2}$, where $\vp=(1+\sqrt5)/2$.
\end{theorem}

\begin{proof}
The $(-2,3,7)$ pretzel knot is listed in the KnotInfo database~\cite{KnotInfo}
as knot $12n242$, with Alexander polynomial $\Delta(t)=t^2+3t+1$.
This value is independently verified by SnapPy~\cite{SnapPy} and
appears in Fintushel--Stern~\cite{FintushelStern1990}, where the
$(-2,3,7)$ pretzel knot plays a central role in instanton homology
calculations for Seifert-fibred homology spheres.

It remains to identify the roots.  The roots of $t^2+3t+1=0$ are
$(-3\pm\sqrt5)/2$.  The golden ratio satisfies $\vp^2=\vp+1$, so
$\vp^2=(3+\sqrt5)/2$ and $\vp^{-2}=(3-\sqrt5)/2$.
Hence the roots are $-\vp^2$ and $-\vp^{-2}$.
This is verified by $\vp^2+\vp^{-2}=(\vp+\vp^{-1})^2-2=(\sqrt5)^2-2=3$
(the coefficient of $t$ in $\Delta$, with sign).
\end{proof}

\begin{corollary}[Mahler measure]\label{cor:mahler}
$M(\Delta)=\vp^2$.
\end{corollary}

\begin{proof}
$M(\Delta)=\prod_i\max(1,|\alpha_i|)=\vp^2\cdot1=\vp^2$,
since $|\vp^2|>1$ and $|\vp^{-2}|<1$.
\end{proof}

\begin{theorem}[Algebraic identity]\label{thm:bridge}
\[
\log M(\Delta) = 2\log\vp = 2\,\mathrm{Reg}(\Q(\sqrt5)).
\]
\end{theorem}

\begin{proof}
From Corollary~\ref{cor:mahler}, $\log M(\Delta)=2\log\vp$.
The regulator of $\Q(\sqrt5)$ is $\log\varepsilon$ where $\varepsilon$
is the fundamental unit.  Since $\vp=(1+\sqrt5)/2$ is the fundamental
unit of $\Q(\sqrt5)$ (the ring of integers is $\Z[\vp]$ and $\vp$ has
norm $-1$, so $|\vp|=\vp>1$ is the smallest unit greater than~1),
$\mathrm{Reg}(\Q(\sqrt5))=\log\vp$.  Hence
$\log M(\Delta)=2\mathrm{Reg}(\Q(\sqrt5))$.
\end{proof}

\begin{remark}
This identity is exact.  The regulator of a real quadratic field
$\Q(\sqrt{d})$ is $\log\varepsilon$, where $\varepsilon>1$ is the
fundamental unit~\cite[Ch.~I, \S12]{Neukirch1999}.  For $\Q(\sqrt5)$,
the ring of integers is $\Z[\vp]$ and $\vp=(1+\sqrt5)/2$ is the
fundamental unit (its norm is $N(\vp)=-1$ and no unit between $1$ and
$\vp$ exists), so $\mathrm{Reg}(\Q(\sqrt5))=\log\vp$ exactly.
\end{remark}

\section{Further arithmetic invariants}
\label{sec:arith}

All computations in this section used \texttt{Sage~10 + SnapPy~3.3}.

\begin{proposition}[Adjoint torsion]\label{prop:torsion}
The adjoint torsion polynomial of \textup{\texttt{m003}} is
\[
\tau_{\mathrm{adj}}(a) = -(a-1)(a^2-5a+1),
\]
with integer coefficients verified to $29$ significant figures.
\end{proposition}

Integrality of the adjoint torsion is a necessary condition for
arithmeticity of a hyperbolic 3-manifold~\cite{MaclachlanReid2003}.

\begin{proposition}[Invariant trace field]\label{prop:tracefield}
The invariant trace field of \textup{\texttt{m003}} is $\Q(\sqrt{-3})$.
Trace field generators satisfy $a^2+5a+7=0$ with discriminant $-3$,
so $\Q(\sqrt{-3})=\Q(\zeta_3)$.
\end{proposition}

\begin{theorem}[Cyclotomic unification]\label{thm:cyclo}
The fields $\Q(\sqrt5)$ and $\Q(\sqrt{-3})$ both embed into the
cyclotomic field $\Q(\zeta_{15})$, with $[\Q(\zeta_{15}):\Q]=8$.
\end{theorem}

\begin{proof}
The subfields of $\Q(\zeta_{15})$ of degree~2 over $\Q$ are determined
by the subgroups of $(\Z/15\Z)^\times\cong\Z/2\times\Z/4$ of index~2.
These subfields are $\Q(\sqrt5)$, $\Q(\sqrt{-3})$, and $\Q(\sqrt{-15})$.
In particular, both $\Q(\sqrt5)$ and $\Q(\sqrt{-3})$ are subfields of
$\Q(\zeta_{15})$.  The degree is $[\Q(\zeta_{15}):\Q]=\varphi(15)=8$.
\end{proof}

\begin{remark}
Theorem~\ref{thm:cyclo} means that the Mahler measure field $\Q(\sqrt5)$
(governing $\log M(\Delta)$) and the invariant trace field $\Q(\sqrt{-3})$
(governing the arithmetic of the holonomy) are unified in a single
cyclotomic extension of degree~8.  The golden ratio $\vp\in\Q(\sqrt5)$
and the trace field generator $\zeta_3\in\Q(\sqrt{-3})$ are thus
arithmetic companions within $\Q(\zeta_{15})$.
\end{remark}

\section{Discussion}
\label{sec:discussion}

The main result of this note is that the Alexander polynomial of the
$(-2,3,7)$ pretzel knot complement has Mahler measure $\vp^2$, and
$\log\vp$ is simultaneously the regulator of $\Q(\sqrt5)$.  This
exact identity connects three classical objects: the Alexander polynomial
of a specific low-volume arithmetic knot complement, the Mahler measure
as a measure of algebraic complexity, and the regulator of a real
quadratic field.

The connection to the figure-eight knot is instructive.  Both
\texttt{m003} and \texttt{m004} are arithmetic hyperbolic knot
complements with the same volume, invariant trace field $\Q(\sqrt{-3})$,
and Chern--Simons invariant $1/4$.  Yet their Alexander polynomials
differ: $\Delta_{4_1}(t)=t-3+t^{-1}$ has Mahler measure $3+2\sqrt2$,
unrelated to $\vp$.  The golden ratio is a feature of \texttt{m003}
specifically, not of the volume class.

The Weeks manifold ($\mathrm{vol}=0.9427$, $H_1=\Z/5\oplus\Z/5$) is
the $(2,1)$ Dehn filling of \texttt{m003}.  Since closed 3-manifolds
do not carry Alexander polynomials, the Mahler measure $\vp^2$ is an
invariant of the cusped manifold only.  The Weeks manifold shares the
arithmetic structure of \texttt{m003} (same cusped ancestor, same
Chern--Simons invariant $1/4$) but the Algebraic Bridge
(Theorem~\ref{thm:bridge}) applies to the cusped manifold alone.

\section{Conclusion}

We have proved: (1)~the Alexander polynomial of the $(-2,3,7)$ pretzel
knot is $t^2+3t+1$ with roots $-\vp^2$, $-\vp^{-2}$;
(2)~its Mahler measure is $\vp^2$;
(3)~$\log M(\Delta)=2\,\mathrm{Reg}(\Q(\sqrt5))$;
(4)~the adjoint torsion has integer coefficients (arithmeticity);
(5)~the invariant trace field $\Q(\sqrt{-3})$ and the Mahler
measure field $\Q(\sqrt5)$ both embed into $\Q(\zeta_{15})$.
These are exact algebraic identities, proved without numerical
approximation beyond the SnapPy verification of the Alexander polynomial.

\bibliographystyle{plain}
\bibliography{gentry-hyperbolic-lattice}

\end{document}
